\begin{titlepage}
\thispagestyle{empty}
\begin{tikzpicture}[remember picture,overlay]
  \fill[AlmanachInk]
    (current page.south west) rectangle (current page.north east);
  \fill[AlmanachTeal]
    (current page.north west) rectangle
    ([yshift=-10mm]current page.north east);
  \draw[AlmanachOchre,line width=1.2pt]
    ([xshift=28mm,yshift=-112mm]current page.north west) --
    ([xshift=-24mm,yshift=-112mm]current page.north east);
  \foreach \x/\y in {
    2.8/5.0,5.1/6.2,7.6/4.7,10.0/6.6,12.4/5.3,14.8/6.8,17.2/5.1}
  {
    \fill[AlmanachTeal!65!white]
      ([xshift=\x cm,yshift=\y cm]current page.south west)
      circle (1.6mm);
  }
  \draw[AlmanachTeal!65!white,line width=0.7pt]
    ([xshift=2.8cm,yshift=5.0cm]current page.south west) --
    ([xshift=5.1cm,yshift=6.2cm]current page.south west) --
    ([xshift=7.6cm,yshift=4.7cm]current page.south west) --
    ([xshift=10.0cm,yshift=6.6cm]current page.south west) --
    ([xshift=12.4cm,yshift=5.3cm]current page.south west) --
    ([xshift=14.8cm,yshift=6.8cm]current page.south west) --
    ([xshift=17.2cm,yshift=5.1cm]current page.south west);
\end{tikzpicture}

\vspace*{31mm}
{\color{white}\sffamily\bfseries\fontsize{38}{43}\selectfont
AI Almanach\par}
\vspace{5mm}
{\color{AlmanachTeal!35!white}\sffamily\fontsize{17}{22}\selectfont
Machine Learning, Data Science, and Intelligent Systems\par}
\vspace{3mm}
{\color{white}\sffamily\large
A prose-first engineering reference\par}

\vfill
\begin{minipage}{0.66\textwidth}
  \color{white!88}
  From mathematical objects to operated systems: each chapter connects
  intuition, formal mechanisms, worked evidence, failure analysis, and
  transfer to practice.
\end{minipage}

\vfill
{\color{white}\sffamily\Large Bernhard Gerlach\par}
\vspace{2mm}
{\color{AlmanachTeal!35!white}\sffamily
Third working edition · \today\par}
\end{titlepage}

\clearpage
\thispagestyle{empty}
\vspace*{\fill}
\small
\begin{tcolorbox}[
  enhanced,
  colback=AlmanachOchre!10,
  colframe=AlmanachOchre,
  boxrule=0pt,
  borderline west={1.4pt}{0pt}{AlmanachOchre},
  arc=1mm,
  left=3mm, right=3mm, top=2.4mm, bottom=2.4mm,
  fonttitle=\small\sffamily\bfseries,
  coltitle=AlmanachOchre!70!black,
  colbacktitle=AlmanachOchre!10,
  titlerule=0pt,
  title={This book was generated by artificial intelligence}]
Every chapter, figure, worked example, and reference list in this Almanach was
written by an AI system, working from a human-defined outline, curriculum, and
editorial standard, and under human direction and review. No part of the text
was written by a human author.

\smallskip
Treat it accordingly. AI systems produce fluent prose whether or not the
underlying claim is correct, and they can misattribute a result, invent a
plausible citation, or carry a subtle numerical error through an otherwise
sound derivation. The safeguards used here---primary sources cited at the
claim they support, worked examples carried through with explicit arithmetic
checks, and derivations stated so they can be re-checked---reduce that risk but
do not eliminate it.

\smallskip
\textbf{Before any consequential use, verify the claim against the cited
primary source.} Where this book and a peer-reviewed paper, a standard, or
official documentation disagree, the primary source is right and this book is
wrong.
\end{tcolorbox}

\vspace{2mm}

\textbf{Purpose.}
This book is an educational engineering reference. It is not a substitute
for professional safety assurance, clinical judgement, financial advice, or
legal review.

\textbf{Evidence.}
Definitions, methods, and empirical claims should be traceable to primary
literature, standards, official documentation, or identified syntheses.
Claims about changing models, software, benchmarks, and law must be checked
again before consequential use.

\textbf{Reproducibility.}
Results depend on data, splits, random seeds, software versions, numerical
precision, hardware, and measurement protocol. A benchmark without its
conditions is a story, not a measurement.

\textbf{Design.}
This edition was rebuilt from a clean document architecture. It uses normal
prose, restrained typographic cues, adjacent explanatory visuals, and
captioned figures and tables rather than a card or poster grid.
\vspace*{\fill}
\clearpage

\chapter*{Abstract}
\addcontentsline{toc}{chapter}{Abstract}

Artificial intelligence is a family of methods for representing problems,
learning from observations, reasoning under uncertainty, optimising choices,
and constructing systems that interact with people and the physical world.
This Almanach develops that family for engineering students and practitioners
whose prior mathematics, programming, and domain experience may differ.

The route begins with mathematical, programming, and data foundations. It
then develops classical machine learning, deep learning, evaluation and
operations, and reinforcement learning. Later chapters connect these
mechanisms to trustworthy AI, perception, language, finance, healthcare,
industry, supply chains, multi-agent systems, commerce, law, visualisation,
venture design, and scientific communication.

Throughout the book, training a model is treated as one link in a larger
chain. Useful systems also require a well-scoped decision, representative
data, meaningful baselines, uncertainty analysis, controlled release,
monitoring, human responsibility, and a credible response when assumptions
fail.

\chapter*{Preface: keeping the reasoning chain visible}
\addcontentsline{toc}{chapter}{Preface}

AI is often difficult because several lessons arrive at once. A learner may
meet a new application, notation, algorithm, and software tool in a single
paragraph. This book separates those layers, then reconnects them through a
repeated engineering loop.

\begin{figure}[htbp]
\centering
\begin{tikzpicture}[
  stage/.style={
    draw=AlmanachTeal,
    fill=AlmanachMist,
    rounded corners=1mm,
    align=center,
    minimum width=27mm,
    minimum height=10mm,
    font=\sffamily\small},
  link/.style={-{Stealth[length=2mm]},AlmanachOchre,thick},
  node distance=8mm and 8mm]
  \node[stage] (question) {Question and\\constraints};
  \node[stage,right=of question] (data) {Data and\\assumptions};
  \node[stage,right=of data] (model) {Model and\\objective};
  \node[stage,below=of model] (evaluate) {Evaluation and\\uncertainty};
  \node[stage,left=of evaluate] (operate) {System and\\operations};
  \node[stage,left=of operate] (revise) {Decision and\\revision};
  \draw[link] (question) -- (data);
  \draw[link] (data) -- (model);
  \draw[link] (model) -- (evaluate);
  \draw[link] (evaluate) -- (operate);
  \draw[link] (operate) -- (revise);
  \draw[link] (revise) -- (question);
\end{tikzpicture}
\caption{The recurring engineering loop connecting questions to evidence.}
\label{fig:frontmatter-engineering-loop}
\end{figure}

The loop matters because a technically elegant model can still fail. Target
leakage, the wrong evaluation unit, an unobserved subgroup, unsafe fallback
behaviour, or silent drift can invalidate the surrounding system even when a
test score looks impressive.

\chapter*{How to use this book}
\addcontentsline{toc}{chapter}{How to use this book}

\section*{Choose a route}

The chapters are dependency ordered, but the book also works as a reference.
Start with the route closest to your goal and return to a prerequisite when a
symbol, assumption, or implementation step is not yet explainable.

\begin{table}[htbp]
\centering
\caption{Suggested learning routes through the AI Almanach.}
\label{tab:frontmatter-learning-routes}
\begin{tblr}{
  width=\textwidth,
  colspec={Q[l,0.19\textwidth]X[l]X[l]},
  row{1}={font=\bfseries,bg=AlmanachMist},
  hlines,
  vlines}
Route & Begin with & Then emphasise \\
Complete first pass & Orientation, mathematics, programming, and data &
Machine learning, deep learning, evaluation, then application chapters \\
ML engineer & Mathematics refresh and data foundations &
ML, deep learning, MLOps, trustworthy AI, and practical labs \\
LLM engineer & Linear algebra, probability, and neural networks &
NLP, practical LLM engineering, evaluation, safety, and operations \\
Data scientist & Statistics, programming, and data analytics &
ML, causal reasoning, visualisation, and deployment \\
Technical leader & Orientation and evaluation &
Trust, law, operations, applications, and evidence limitations \\
\end{tblr}
\end{table}

\section*{A repeated learning sequence}

The sequence below makes the reasoning visible and exposes missing evidence.
It combines constructive alignment, worked-example support, retrieval
practice, and spatial contiguity between words and meaningful graphics
\cite{biggsEnhancingTeachingConstructive1996,
swellerCognitiveLoadTheory2019,dunloskyImprovingStudentsLearning2013,
roedigerTestEnhancedLearning2006,mayerUsingMultimediaElearning2017}.

\begin{enumerate}
  \item Start with the decision or question that gives the method a purpose.
  \item State prerequisites, notation, units, and operating assumptions.
  \item Build a plain-language mental model and identify where it breaks.
  \item Follow the formal mechanism through a small, worked example.
  \item Connect the mechanism to implementation and systems constraints.
  \item Predict a failure, then design a test that can distinguish it.
  \item Retrieve the reasoning without looking and transfer it to a new case.
\end{enumerate}

\begin{table}[htbp]
\centering
\caption{Evidence-informed learning patterns used in every chapter.}
\label{tab:frontmatter-didactic-patterns}
\begin{tblr}{
  width=\textwidth,
  colspec={Q[l,0.22\textwidth]X[l]X[l]},
  row{1}={font=\bfseries,bg=AlmanachMist},
  hlines,
  vlines}
Pattern & Purpose & Reader action \\
Chapter roadmap & Expose dependencies before detail &
Predict the four-stage chain, then redraw it after reading \\
Worked example & Make the reasoning path observable &
Explain each step, then complete a changed case \\
Failure analysis & Build diagnostic knowledge &
Locate the violated assumption and propose a discriminating test \\
Retrieval prompt & Strengthen later recall and reveal gaps &
Answer without looking, check, correct, and revisit after a delay \\
Transfer task & Prevent example-bound knowledge &
Change the domain or constraints and adapt the method \\
Adjacent visual & Reduce split attention for spatial ideas &
Trace the figure while explaining the corresponding prose \\
\end{tblr}
\end{table}

\section*{Evidence and notation}

Use four evidence levels: primary or normative evidence; authoritative
synthesis; official operational documentation; and illustrative examples.
The final category can explain an idea but cannot establish a general claim.
For empirical results, record the dataset, split, metric, uncertainty,
hardware, and software configuration. For equations, check dimensions,
boundary cases, sign conventions, and at least one numerical example.

\begin{table}[htbp]
\centering
\caption{Core notation used across chapters.}
\label{tab:frontmatter-core-notation}
\begin{tblr}{
  width=\textwidth,
  colspec={Q[c,0.14\textwidth]X[l]X[l]},
  row{1}={font=\bfseries,bg=AlmanachMist},
  hlines,
  vlines}
Symbol & Meaning & Typical shape or range \\
$\mathbf{x}$ & Feature vector & $\mathbf{x}\in\mathbb{R}^{d}$ \\
$X$ & Design matrix or token representations &
$X\in\mathbb{R}^{n\times d}$ \\
$y,\hat y$ & Observed target and model prediction &
Scalar, class, sequence, or structured object \\
$\theta$ & Trainable parameters & Vector or collection of tensors \\
$\mathcal{L}$ & Loss or empirical objective & Usually $\mathbb{R}$ \\
$p(y\mid x)$ & Conditional probability &
$[0,1]$, summing or integrating to one \\
$\nabla_{\theta}\mathcal{L}$ & Gradient with respect to parameters &
The same shape as $\theta$ \\
$\mathbb{E}[X]$ & Expectation of a random variable & The units of $X$ \\
\end{tblr}
\end{table}

\begin{checkpoint}
Choose a model you already know. Name its input, target, parameters,
objective, evaluation unit, one assumption, one failure mode, and the person
or process responsible for the final decision. Every blank is a useful
reading goal.
\end{checkpoint}
